% SOURCE_AUTHORITY: sga5_fr_workpass.tex, lines 3961--4013 (LF-normalized unit).
% SOURCE_SHA256: 791F4EFFC5E02832D5D77ED03518C8156D6F07E4C8238B03545DB93D883FBB28
\subsection*{6.9}

Sean $X$ (resp. $Y$) un $S$-esquema propio y liso de dimensión relativa $m$ (resp. $n$), $Z=X\times_S Y$, y sean $a:A\hookrightarrow Z$, $b:B\hookrightarrow Z$ inmersiones regulares tales que $a_2:A\to Y$ y $b_1:B\to X$ sean finitos y planos (denotamos por $p_1:Z\to X$, $p_2:Z\to Y$ las proyecciones, y $a_i=p_i a$, $b_i=p_i b$), $c:C=A\times_Z B\to Z$, $L\in\ob D(X)_{\parf}$ y $M\in\ob D(Y)_{\parf}$ (el índice parf significa ``perfecto en sentido absoluto'', lo que equivale también a perfecto relativo a $S$, pues $X$ e $Y$ son lisos). Se supone que $C$ es finito y plano sobre $S$. Las hipótesis impuestas implican, por tanto, que $c$ es una inmersión regular y que $a$ y $b$ son Tor-independientes.

Se tiene $p_1^*\Omega^1_{X/S}=\Omega^1_{Z/Y}$, $p_2^*\Omega^1_{Y/S}=\Omega^1_{Z/X}$, y una descomposición
\[
\Omega^1_{Z/S}=\Omega^1_{Z/Y}\oplus\Omega^1_{Z/X},
\]
de donde, para todo entero $r$, se obtiene una descomposición
\begin{equation}
\tag{6.9.0}
\Omega^r_{Z/S}=\bigoplus_{i+j=r}\Omega^i_{Z/Y}\otimes\Omega^j_{Z/X}.
\end{equation}
Por otra parte, según 6.1.2, se tiene $p_2^!M=p_2^*M\otimes\Omega^m_{Z/Y}[m]$, y, por tanto,
\[
H^0(Z,a_*a^!\uRHom(p_1^*L,p_2^!M))=\Ext^m_{\OO_Z}(a_1^*L,a_2^*M\otimes\Omega^m_{Z/Y}),
\]
de donde, por 6.4.4 y 6.8.2, se obtiene un isomorfismo
\begin{equation}
\tag{6.9.1}
\Hom(a_1^*L,a_2^!M)=H^0(A,\uRHom(a_1^*L,a_2^*M)\otimes\Lambda^m N_{A/Z}^{\vee}\otimes\Omega^m_{Z/Y}).
\end{equation}
Del mismo modo, se tiene un isomorfismo
\begin{equation}
\tag{6.9.2}
\Hom(b_2^*M,b_1^!L)=H^0(B,\uRHom(b_2^*M,b_1^*L)\otimes\Lambda^n N_{B/Z}^{\vee}\otimes\Omega^n_{Z/X}).
\end{equation}
Sean $u\in\Hom(a_1^*L,a_2^*M)$ y $v\in\Hom(b_2^*M,b_1^*L)$; denotemos por
\begin{equation}
\tag{6.9.3}
\langle u,v\rangle\in H^0(C,\OO_C)
\end{equation}
la sección $\operatorname{Tr}(u_Cv_C)=\operatorname{Tr}(v_Cu_C)$, donde $u_C\in\Hom(c_1^*L,c_2^*M)$ y $v_C\in\Hom(c_2^*M,c_1^*L)$ están inducidos por $u$ y $v$. Consideremos la correspondencia cohomológica $u^!\in\Hom(a_1^*L,a_2^!M)$ (resp. $v^!\in\Hom(b_2^*M,b_1^!L)$) deducida de $u$ (resp. $v$) componiendo con el homomorfismo de traza 6.8.7 $a_{2*}a_2^*M\to M$ (resp. $b_{1*}b_1^*L\to L$), o, lo que es equivalente (6.8.6), tomando el producto con $\cl_{Z/Y}(A)$ (resp. $\cl_{Z/X}(B)$). Observemos de paso que $\cl_{Z/Y}(A)$ (resp. $\cl_{Z/X}(B)$) es la componente en $H^0(A,\Lambda^mN_{A/Z}^{\vee}\otimes\Omega^m_{Z/Y})$ (resp. $H^0(B,\Lambda^nN_{B/Z}^{\vee}\otimes\Omega^n_{Z/X})$) de $\cl_{Z/S}(A)$ (resp. $\cl_{Z/S}(B)$) definida por la proyección $\Omega^m_{Z/S}\to\Omega^m_{Z/Y}$ (resp. $\Omega^n_{Z/S}\to\Omega^n_{Z/X}$) de 6.9.0. Se comprueba fácilmente que el producto cup
\[
\langle u^!,v^!\rangle\in\Ext^{m+n}_{\OO_Z}(\OO_C,\Omega^{m+n}_{Z/S})=H^0(C,\Lambda^{m+n}N_{C/Z}^{\vee}\otimes\Omega^{m+n}_{Z/S})
\]
definido en 6.6.3 viene dado por
\begin{equation}
\tag{6.9.4}
\langle u^!,v^!\rangle=\langle u,v\rangle\cl_{Z/Y}(A)\cl_{Z/X}(B).
\end{equation}
Por otra parte, se comprueba que el homomorfismo 6.7.3
\[
\Ext^{m+n}_{\OO_Z}(\OO_C,\Omega^{m+n}_{Z/S})\longrightarrow H^0(S,\OO_S)
\]
es la compuesta de la flecha canónica $\Ext^{m+n}_{\OO_Z}(\OO_C,\Omega^{m+n}_{Z/S})\to H^{m+n}(Z,\Omega^{m+n}_{Z/S})$ y del morfismo de traza $\operatorname{Tr}_{Z/S}:H^{m+n}(Z,\Omega^{m+n}_{Z/S})\to H^0(S,\OO_S)$. La fórmula de Lefschetz--Verdier 6.7.2 proporciona, por tanto, la relación
\begin{equation}
\tag{6.9.5}
\langle (u^!)_*,(v^!)_*\rangle=\operatorname{Tr}_{Z/S}(\langle u,v\rangle\cl_{Z/Y}(A)\cl_{Z/X}(B)).
\end{equation}
Gracias a 6.8.4, el segundo miembro puede reescribirse como un residuo; en resumen, hemos obtenido:


